% Neural-network architecture section — % Neural-network architecture section — % Neural-network architecture section — % Neural-network architecture section — \input{nn_architectures} in main paper

\subsection{Neural-network architectures}
\label{sec:nn_arch}

Three binary classifiers are trained to discriminate VBF~$H\!\to\!\text{inv}$ signal
from the QCD multijet background.
Each classifier is a fully-connected multilayer perceptron (MLP) implemented in
PyTorch~\cite{pytorch}.
The three networks are identical in every architectural and training detail;
the \emph{only} difference between them is the set of input features, as
described below.

\subsubsection*{Network architecture}

Each MLP maps an input vector $\mathbf{x}\in\mathbb{R}^{N_\text{in}}$ to a
single scalar output $\hat{y}\in(0,1)$ through three hidden layers of widths
$128$, $64$, and $32$, followed by a sigmoid activation:
%
\begin{equation}
  \hat{y} = \sigma\!\left(
    W_4\,g\!\left(W_3\,g\!\left(W_2\,g\!\left(W_1\,\mathbf{x}+b_1\right)+b_2\right)+b_3\right)+b_4
  \right),
\end{equation}
%
where $g$ denotes the composed activation block (BatchNorm $\to$ ReLU $\to$
Dropout) applied after each hidden linear layer, and $\sigma$ is the sigmoid
function.
Concretely, each hidden layer consists of:
%
\begin{enumerate}
  \item a linear transformation,
  \item batch normalisation~\cite{batchnorm},
  \item a ReLU activation,
  \item dropout with probability $p=0.3$.
\end{enumerate}
%
The output layer is a single linear unit followed by a sigmoid, yielding a
signal probability score.
All input features are standardised to zero mean and unit variance using
statistics computed on the training set only.

\subsubsection*{Input features}

The three models differ solely in which features are passed as input:
%
\begin{table}[h]
  \centering
  \caption{Input features for each of the three classifiers.
           $m_{jj}$: dijet invariant mass;
           $|\Delta y_{jj}|$: absolute rapidity separation of the two leading jets;
           $|\Delta\phi_{jj}|$: absolute azimuthal-angle separation;
           $p_{T,jj}$: dijet transverse momentum;
           $|\text{SPVA}[j_{1,2}]|$: magnitude of the signed pull-vector angle
           for the leading (subleading) jet;
           $\text{PVM}[j_{1,2}]$: pull-vector magnitude for the
           leading (subleading) jet.}
  \label{tab:nn_features}
  \renewcommand{\arraystretch}{1.25}
  \begin{tabular}{clcc}
    \hline\hline
    Model & Description & $N_\text{in}$ & Input features \\
    \hline
    A & Kinematic & 4 &
      $m_{jj},\;|\Delta y_{jj}|,\;|\Delta\phi_{jj}|,\;p_{T,jj}$ \\
    B & Colour / pull & 4 &
      $|\text{SPVA}[j_1]|,\;|\text{SPVA}[j_2]|,\;\text{PVM}[j_1],\;\text{PVM}[j_2]$ \\
    C & Combined & 8 &
      all features from models A and B \\
    \hline\hline
  \end{tabular}
\end{table}

Model~A uses only standard VBF kinematic discriminants.
Model~B uses only jet colour-flow observables (pull vector angle and magnitude),
which are sensitive to the colour singlet nature of the VBF topology.
Model~C combines both sets and therefore has a wider input layer
($N_\text{in}=8$) while keeping the hidden-layer widths unchanged.

\subsubsection*{Training details}

All three networks are trained with the same procedure:
binary cross-entropy loss, the Adam optimiser~\cite{adam} with initial learning
rate $\eta=10^{-3}$, and a \textsc{ReduceLROnPlateau} scheduler that halves
$\eta$ after five epochs without improvement in the validation loss.
Training proceeds for up to 200 epochs with early stopping: training is halted
if the validation loss does not improve for 20 consecutive epochs, and the
checkpoint with the lowest validation loss is retained.
The mini-batch size is 4096 events.
The full dataset is split 60\,/\,20\,/\,20 into training, validation, and
held-out inference sets using a fixed random seed.
Events are required to satisfy $m_{jj}>400\;\text{GeV}$.
 in main paper

\subsection{Neural-network architectures}
\label{sec:nn_arch}

Three binary classifiers are trained to discriminate VBF~$H\!\to\!\text{inv}$ signal
from the QCD multijet background.
Each classifier is a fully-connected multilayer perceptron (MLP) implemented in
PyTorch~\cite{pytorch}.
The three networks are identical in every architectural and training detail;
the \emph{only} difference between them is the set of input features, as
described below.

\subsubsection*{Network architecture}

Each MLP maps an input vector $\mathbf{x}\in\mathbb{R}^{N_\text{in}}$ to a
single scalar output $\hat{y}\in(0,1)$ through three hidden layers of widths
$128$, $64$, and $32$, followed by a sigmoid activation:
%
\begin{equation}
  \hat{y} = \sigma\!\left(
    W_4\,g\!\left(W_3\,g\!\left(W_2\,g\!\left(W_1\,\mathbf{x}+b_1\right)+b_2\right)+b_3\right)+b_4
  \right),
\end{equation}
%
where $g$ denotes the composed activation block (BatchNorm $\to$ ReLU $\to$
Dropout) applied after each hidden linear layer, and $\sigma$ is the sigmoid
function.
Concretely, each hidden layer consists of:
%
\begin{enumerate}
  \item a linear transformation,
  \item batch normalisation~\cite{batchnorm},
  \item a ReLU activation,
  \item dropout with probability $p=0.3$.
\end{enumerate}
%
The output layer is a single linear unit followed by a sigmoid, yielding a
signal probability score.
All input features are standardised to zero mean and unit variance using
statistics computed on the training set only.

\subsubsection*{Input features}

The three models differ solely in which features are passed as input:
%
\begin{table}[h]
  \centering
  \caption{Input features for each of the three classifiers.
           $m_{jj}$: dijet invariant mass;
           $|\Delta y_{jj}|$: absolute rapidity separation of the two leading jets;
           $|\Delta\phi_{jj}|$: absolute azimuthal-angle separation;
           $p_{T,jj}$: dijet transverse momentum;
           $|\text{SPVA}[j_{1,2}]|$: magnitude of the signed pull-vector angle
           for the leading (subleading) jet;
           $\text{PVM}[j_{1,2}]$: pull-vector magnitude for the
           leading (subleading) jet.}
  \label{tab:nn_features}
  \renewcommand{\arraystretch}{1.25}
  \begin{tabular}{clcc}
    \hline\hline
    Model & Description & $N_\text{in}$ & Input features \\
    \hline
    A & Kinematic & 4 &
      $m_{jj},\;|\Delta y_{jj}|,\;|\Delta\phi_{jj}|,\;p_{T,jj}$ \\
    B & Colour / pull & 4 &
      $|\text{SPVA}[j_1]|,\;|\text{SPVA}[j_2]|,\;\text{PVM}[j_1],\;\text{PVM}[j_2]$ \\
    C & Combined & 8 &
      all features from models A and B \\
    \hline\hline
  \end{tabular}
\end{table}

Model~A uses only standard VBF kinematic discriminants.
Model~B uses only jet colour-flow observables (pull vector angle and magnitude),
which are sensitive to the colour singlet nature of the VBF topology.
Model~C combines both sets and therefore has a wider input layer
($N_\text{in}=8$) while keeping the hidden-layer widths unchanged.

\subsubsection*{Training details}

All three networks are trained with the same procedure:
binary cross-entropy loss, the Adam optimiser~\cite{adam} with initial learning
rate $\eta=10^{-3}$, and a \textsc{ReduceLROnPlateau} scheduler that halves
$\eta$ after five epochs without improvement in the validation loss.
Training proceeds for up to 200 epochs with early stopping: training is halted
if the validation loss does not improve for 20 consecutive epochs, and the
checkpoint with the lowest validation loss is retained.
The mini-batch size is 4096 events.
The full dataset is split 60\,/\,20\,/\,20 into training, validation, and
held-out inference sets using a fixed random seed.
Events are required to satisfy $m_{jj}>400\;\text{GeV}$.
 in main paper

\subsection{Neural-network architectures}
\label{sec:nn_arch}

Three binary classifiers are trained to discriminate VBF~$H\!\to\!\text{inv}$ signal
from the QCD multijet background.
Each classifier is a fully-connected multilayer perceptron (MLP) implemented in
PyTorch~\cite{pytorch}.
The three networks are identical in every architectural and training detail;
the \emph{only} difference between them is the set of input features, as
described below.

\subsubsection*{Network architecture}

Each MLP maps an input vector $\mathbf{x}\in\mathbb{R}^{N_\text{in}}$ to a
single scalar output $\hat{y}\in(0,1)$ through three hidden layers of widths
$128$, $64$, and $32$, followed by a sigmoid activation:
%
\begin{equation}
  \hat{y} = \sigma\!\left(
    W_4\,g\!\left(W_3\,g\!\left(W_2\,g\!\left(W_1\,\mathbf{x}+b_1\right)+b_2\right)+b_3\right)+b_4
  \right),
\end{equation}
%
where $g$ denotes the composed activation block (BatchNorm $\to$ ReLU $\to$
Dropout) applied after each hidden linear layer, and $\sigma$ is the sigmoid
function.
Concretely, each hidden layer consists of:
%
\begin{enumerate}
  \item a linear transformation,
  \item batch normalisation~\cite{batchnorm},
  \item a ReLU activation,
  \item dropout with probability $p=0.3$.
\end{enumerate}
%
The output layer is a single linear unit followed by a sigmoid, yielding a
signal probability score.
All input features are standardised to zero mean and unit variance using
statistics computed on the training set only.

\subsubsection*{Input features}

The three models differ solely in which features are passed as input:
%
\begin{table}[h]
  \centering
  \caption{Input features for each of the three classifiers.
           $m_{jj}$: dijet invariant mass;
           $|\Delta y_{jj}|$: absolute rapidity separation of the two leading jets;
           $|\Delta\phi_{jj}|$: absolute azimuthal-angle separation;
           $p_{T,jj}$: dijet transverse momentum;
           $|\text{SPVA}[j_{1,2}]|$: magnitude of the signed pull-vector angle
           for the leading (subleading) jet;
           $\text{PVM}[j_{1,2}]$: pull-vector magnitude for the
           leading (subleading) jet.}
  \label{tab:nn_features}
  \renewcommand{\arraystretch}{1.25}
  \begin{tabular}{clcc}
    \hline\hline
    Model & Description & $N_\text{in}$ & Input features \\
    \hline
    A & Kinematic & 4 &
      $m_{jj},\;|\Delta y_{jj}|,\;|\Delta\phi_{jj}|,\;p_{T,jj}$ \\
    B & Colour / pull & 4 &
      $|\text{SPVA}[j_1]|,\;|\text{SPVA}[j_2]|,\;\text{PVM}[j_1],\;\text{PVM}[j_2]$ \\
    C & Combined & 8 &
      all features from models A and B \\
    \hline\hline
  \end{tabular}
\end{table}

Model~A uses only standard VBF kinematic discriminants.
Model~B uses only jet colour-flow observables (pull vector angle and magnitude),
which are sensitive to the colour singlet nature of the VBF topology.
Model~C combines both sets and therefore has a wider input layer
($N_\text{in}=8$) while keeping the hidden-layer widths unchanged.

\subsubsection*{Training details}

All three networks are trained with the same procedure:
binary cross-entropy loss, the Adam optimiser~\cite{adam} with initial learning
rate $\eta=10^{-3}$, and a \textsc{ReduceLROnPlateau} scheduler that halves
$\eta$ after five epochs without improvement in the validation loss.
Training proceeds for up to 200 epochs with early stopping: training is halted
if the validation loss does not improve for 20 consecutive epochs, and the
checkpoint with the lowest validation loss is retained.
The mini-batch size is 4096 events.
The full dataset is split 60\,/\,20\,/\,20 into training, validation, and
held-out inference sets using a fixed random seed.
Events are required to satisfy $m_{jj}>400\;\text{GeV}$.
 in main paper

\subsection{Neural-network architectures}
\label{sec:nn_arch}

Three binary classifiers are trained to discriminate VBF~$H\!\to\!\text{inv}$ signal
from the QCD multijet background.
Each classifier is a fully-connected multilayer perceptron (MLP) implemented in
PyTorch~\cite{pytorch}.
The three networks are identical in every architectural and training detail;
the \emph{only} difference between them is the set of input features, as
described below.

\subsubsection*{Network architecture}

Each MLP maps an input vector $\mathbf{x}\in\mathbb{R}^{N_\text{in}}$ to a
single scalar output $\hat{y}\in(0,1)$ through three hidden layers of widths
$128$, $64$, and $32$, followed by a sigmoid activation:
%
\begin{equation}
  \hat{y} = \sigma\!\left(
    W_4\,g\!\left(W_3\,g\!\left(W_2\,g\!\left(W_1\,\mathbf{x}+b_1\right)+b_2\right)+b_3\right)+b_4
  \right),
\end{equation}
%
where $g$ denotes the composed activation block (BatchNorm $\to$ ReLU $\to$
Dropout) applied after each hidden linear layer, and $\sigma$ is the sigmoid
function.
Concretely, each hidden layer consists of:
%
\begin{enumerate}
  \item a linear transformation,
  \item batch normalisation~\cite{batchnorm},
  \item a ReLU activation,
  \item dropout with probability $p=0.3$.
\end{enumerate}
%
The output layer is a single linear unit followed by a sigmoid, yielding a
signal probability score.
All input features are standardised to zero mean and unit variance using
statistics computed on the training set only.

\subsubsection*{Input features}

The three models differ solely in which features are passed as input:
%
\begin{table}[h]
  \centering
  \caption{Input features for each of the three classifiers.
           $m_{jj}$: dijet invariant mass;
           $|\Delta y_{jj}|$: absolute rapidity separation of the two leading jets;
           $|\Delta\phi_{jj}|$: absolute azimuthal-angle separation;
           $p_{T,jj}$: dijet transverse momentum;
           $|\text{SPVA}[j_{1,2}]|$: magnitude of the signed pull-vector angle
           for the leading (subleading) jet;
           $\text{PVM}[j_{1,2}]$: pull-vector magnitude for the
           leading (subleading) jet.}
  \label{tab:nn_features}
  \renewcommand{\arraystretch}{1.25}
  \begin{tabular}{clcc}
    \hline\hline
    Model & Description & $N_\text{in}$ & Input features \\
    \hline
    A & Kinematic & 4 &
      $m_{jj},\;|\Delta y_{jj}|,\;|\Delta\phi_{jj}|,\;p_{T,jj}$ \\
    B & Colour / pull & 4 &
      $|\text{SPVA}[j_1]|,\;|\text{SPVA}[j_2]|,\;\text{PVM}[j_1],\;\text{PVM}[j_2]$ \\
    C & Combined & 8 &
      all features from models A and B \\
    \hline\hline
  \end{tabular}
\end{table}

Model~A uses only standard VBF kinematic discriminants.
Model~B uses only jet colour-flow observables (pull vector angle and magnitude),
which are sensitive to the colour singlet nature of the VBF topology.
Model~C combines both sets and therefore has a wider input layer
($N_\text{in}=8$) while keeping the hidden-layer widths unchanged.

\subsubsection*{Training details}

All three networks are trained with the same procedure:
binary cross-entropy loss, the Adam optimiser~\cite{adam} with initial learning
rate $\eta=10^{-3}$, and a \textsc{ReduceLROnPlateau} scheduler that halves
$\eta$ after five epochs without improvement in the validation loss.
Training proceeds for up to 200 epochs with early stopping: training is halted
if the validation loss does not improve for 20 consecutive epochs, and the
checkpoint with the lowest validation loss is retained.
The mini-batch size is 4096 events.
The full dataset is split 60\,/\,20\,/\,20 into training, validation, and
held-out inference sets using a fixed random seed.
Events are required to satisfy $m_{jj}>400\;\text{GeV}$.
